\documentclass[11pt]{article}
\usepackage{amsmath,amssymb,amsthm,hyperref}
\newtheorem{proposition}{Proposition}
\title{A Degree Audit for the Discrete-Input Sigmoid Argument}
\author{Research checkpoint, branch ar/027ac8030b43}
\date{November 24, 2026}
\begin{document}
\maketitle

\paragraph{Scope.}
This note identifies an explicit false step in the printed proof of
Theorem 8.11 of Anthony and Bartlett, \emph{Neural Network Learning:
Theoretical Foundations}, printed pages 122--123 (PDF pages 136--137 in
\href{https://openeclass.panteion.gr/modules/document/file.php/PMS152/LEARNING/Anthony\%20Martin\%2C\%20Bartlett\%20Peter\%20L.\%20\%282009\%29\%20--\%20Neural\%20Network\%20Learning_\%20Theoretical\%20Foundations.pdf}{the inspected copy}).
It does not refute the theorem's VC-dimension conclusion, which could
conceivably have a different proof. It does not establish an obstruction
to the original xorformer benchmark. No model definition is changed.
The counterexample and degree argument here are written proofs, not
Lean declarations.

\section{The first failing assertion}

The printed theorem gives an upper bound $2W\log_2(60ND)$ for two-layer
sigmoid networks on $\{-D,\ldots,D\}^n$, with at most $N$ inputs per
hidden unit and $W$ total parameters. Its proof transforms first-layer
weights and thresholds by exponentiation. One hidden unit then becomes
a rational function whose numerator and denominator have degrees at most
$ND$ and $2ND+1$, respectively.

The next step says that an affine combination of $k$ such units has the
same sign as a polynomial of degree at most $3ND+2$ in all transformed
parameters. This assertion fails even for $N=D=1$ and seven units.
The error is not resolved by distinguishing ordinary degree from degree
in each parameter block: Theorem 8.3 uses ordinary total degree, and the
following affine-line restriction rules out any polynomial of the claimed
ordinary degree, regardless of how it is obtained.

\begin{proposition}
For a fixed input of a seven-unit, one-input sigmoid network, the sign
function in the transformed parameters need not admit a polynomial sign
representation of total degree at most five on the positive parameter
orthant.
\end{proposition}

\begin{proof}
Fix the input $x=0$. Along an affine line in the transformed parameter
space, set the threshold parameters to $a_{0,i}(t)=t+i-1$, for
$i=1,\ldots,7$ and $t>0$, and set the transformed input weights to one.
All transformed first-layer parameters are strictly positive. The hidden
outputs are $1/(t+i)$. Choose output bias zero and fixed output weights
\[
(c_1,\ldots,c_7)=(7,-168,1260,-4200,6930,-5544,1716).
\]
Their sum is exactly
\[
F(t)=\sum_{i=1}^7\frac{c_i}{t+i}
 =\frac{\prod_{j=1}^6(t-j)}{\prod_{i=1}^7(t+i)}.
\tag{1}
\]
For example, this identity follows from partial fractions, since
$c_i=\prod_{j=1}^6(-i-j)/\prod_{\ell\ne i}(-i+\ell)$.
The denominator is positive for every $t>0$, and the numerator has the
six distinct simple positive roots $1,\ldots,6$. Thus the seven values
$F(1/2),F(3/2),\ldots,F(13/2)$ have alternating, nonzero signs.

A polynomial in all network parameters of total degree at most five
restricts on this affine parameter line to a univariate polynomial of
degree at most five. If it represented the same strict signs, the
intermediate value theorem would give a distinct root between every
consecutive pair of the seven displayed points. Six distinct roots are
impossible unless the polynomial is zero, which also cannot represent
the nonzero signs. Here $3ND+2=5$, proving the assertion.
\end{proof}

For reproducibility, the seven exact rational values in the proof are
\[
\frac2{195},\quad-\frac2{12155},\quad\frac2{138567},\quad
-\frac{10}{2909907},\quad\frac2{1062347},\quad
-\frac6{2414425},\quad\frac{22}{1671525}.
\]
They were independently evaluated with Python's exact rational arithmetic;
the displayed partial-fraction identity supplies the general proof.

\section{What straightforward clearing actually gives}

Write each hidden output as $P_i/Q_i$, with $Q_i>0$, and let the output
score be $\beta+\sum_i c_i P_i/Q_i$. Clearing all positive denominators
gives
\[
\beta\prod_iQ_i+\sum_i c_iP_i\prod_{\ell\ne i}Q_\ell.
\]
If each hidden unit has at most $N$ inputs, this polynomial has degree at
most $1+k(2ND+1)$. If the individual fan-ins are $N_i$, the sharper bound
is $1+2D\sum_iN_i+k$, hence $O(DW)$ when $D\geq1$. Applying the ordinary
polynomial sign-count argument at this degree yields $O(W\log(WD))$,
with the total parameter count inside the logarithm. It does not yield
the claimed replacement of $W$ by the maximum hidden fan-in $N$.

The book's bibliographical notes, printed page 129, explicitly describe
Theorems 8.11 and 8.12 as a fan-in refinement of Bartlett and Williamson
(1996), DOI \href{https://doi.org/10.1162/neco.1996.8.3.625}
{10.1162/neco.1996.8.3.625}. The present audit did not retrieve that paper's
full text from an accessible primary endpoint, so it does not claim to
have audited its proof. Targeted searches of the authors' publication
pages and for a book erratum did not locate an erratum addressing this
step as of November 24, 2026. Failure to find an erratum is not evidence
that none exists.

\section{Implication for the xorformer investigation}

A sum of $H$ positive-denominator affine ratios has the same elementary
clearing issue: ordinary total degree grows with $H$. The bounded degree
of each individual head does not justify dropping $\log H$ from the
polynomial sign-pattern count. A stronger multiblock or structural
count would require an independent proof. The book argument therefore
cannot be used here to rule out a fixed-power or unbounded-factor
improvement in worst-case head complexity.
\end{document}
